
\begin{proposition}[Credited re-derivation]\label{cdp:st:D.1.1}
For an action $\rho:Q\to\operatorname{Aut}(K)$ and the corresponding
semidirect product $G=K\rtimes_\rho Q$, put
\[
(K^{\rm ab})_Q=K^{\rm ab}/
\langle\rho(q)\bar k-\bar k:q\in Q,\ \bar k\in K^{\rm ab}\rangle.
\]
There is an isomorphism
\[
G^{\rm ab}\longrightarrow Q^{\rm ab}\oplus(K^{\rm ab})_Q,
\qquad [(k,q)]\longmapsto([q],[k]).
\]
For a split extension without a specified splitting, the associated
short exact sequence is canonical and its splitting need not be.
\end{proposition}

This semidirect-product abelianization is \citet[Lemma~10.7, p.~90]{AC2026}. We retain the complete re-derivation below, including the dependence of the direct-sum splitting on the chosen splitting of the extension.

\begin{proof}
The product is $(k,q)(k',q')=(k\rho(q)(k'),qq')$.
In the displayed quotient $[\rho(q)(k')]=[k']$, so the displayed map
is a surjective homomorphism to an abelian group. Its kernel contains
$[K,K]$, $[Q,Q]$ and all elements $\rho(q)(k)k^{-1}$.
These elements are themselves commutators in $G$. Conversely their
normal closure makes the images of $K$ and $Q$ abelian and makes
those images commute, so the quotient by that normal closure is
abelian. The kernel is exactly $[G,G]$. The inverse homomorphism
sends the two summands to the images of $(1,q)$ and $(k,1)$.
The induced maps in
$0\to(K^{\rm ab})_Q\to G^{\rm ab}\to Q^{\rm ab}\to0$
come from inclusion and projection; choosing the original group
section supplies its splitting.

For a concrete failure of the uncorrected direct sum, take
$K=Q=\Z$ and let the generator of $Q$ act on $K$ by $k\mapsto-k$.
Then the coinvariant quotient is $\Z/2\Z$, and the displayed map
gives $G^{\rm ab}=\Z\oplus\Z/2\Z$, whereas
$Q^{\rm ab}\oplus K^{\rm ab}=\Z^2$.
\end{proof}


